\documentclass[12pt]{article}

\usepackage[margin=1in]{geometry}
\usepackage[T1]{fontenc}
\usepackage{lmodern}
\usepackage{microtype}
\usepackage{amsmath,amssymb,mathtools}
\usepackage{booktabs}
\usepackage{setspace}
\usepackage{hyperref}
\usepackage{xurl}

\hypersetup{
  colorlinks=true,
  linkcolor=black,
  citecolor=black,
  urlcolor=blue
}

\onehalfspacing
\widowpenalty=10000
\clubpenalty=10000

\title{A Taylor-Series Expansion of Relativistic Energy:\\
Rest Energy, Kinetic Energy, and the Newtonian Limit}
\author{Light-Meter Physics}
\date{}

\begin{document}

\maketitle

\begin{abstract}
The relativistic energy of a free massive particle can be written as a simple function of its speed. Expanding that function about zero speed makes the relation between total energy, rest energy, and the Newtonian kinetic-energy approximation explicit. This paper derives the Taylor coefficients of the dimensionless function $(1-x)^{-1/2}$, restores physical units, and identifies the domain in which the infinite series converges. The resulting kinetic-energy series begins with $mv^2/2$ and continues with positive relativistic corrections. Numerical examples show why a short truncation is reliable at low speed but can become inaccurate as the speed approaches the speed of light. The relativistic energy relation is taken as the starting point; the purpose here is to explain its series expansion, not to derive special relativity from first principles.
\end{abstract}

\noindent\textbf{Keywords:} special relativity, relativistic energy, Taylor series, binomial expansion, kinetic energy

\section{Introduction}

The familiar Newtonian expression for kinetic energy,
\begin{equation}
  K_{\mathrm{N}}=\frac{1}{2}mv^2,
  \label{eq:newtonian-kinetic-energy}
\end{equation}
is an excellent approximation when an object's speed is small compared with the speed of light. Special relativity replaces this expression with a formula that remains valid at every speed below the speed of light. Expanding the relativistic formula at low speed shows exactly how the Newtonian result appears and how the first corrections depend on higher powers of $v/c$.

The goal of this note is pedagogical. It presents one derivation in a sequence that keeps the physical quantities, dimensionless mathematics, and approximation limits distinct. First, the exact total energy is stated. Next, a dimensionless variable is chosen so that the binomial series can be applied. The Taylor coefficients are then calculated, the physical units are restored, and the resulting terms are interpreted as rest and kinetic energy. Finally, the convergence domain is distinguished from the accuracy of a finite truncation.

This manuscript expands the seven-panel carousel derivation into a paper-length explanation \cite{carousel}. The Taylor expansion is a standard consequence of the binomial series \cite{DLMF}, and no new physical law is proposed. In particular, the calculation does not establish the Lorentz factor or the postulates of special relativity; it starts from the relativistic energy of a free massive particle and analyzes that expression. Einstein's 1905 treatment of electron kinetic energy provides historical context, but the calculation below is a modern series expansion rather than a reproduction of that argument \cite{Einstein1905}.

This is the velocity-based treatment in the numbered Papers series. A companion paper studies low- and high-momentum expansions of the same energy--momentum relation \cite{EnergyMomentum}, while another examines the proper-time rate and its exact relation to kinetic energy \cite{ProperTime}. Together, these papers use $m$ for the particle's invariant mass.

\section{Exact energy and notation}

Consider a free particle with invariant mass $m>0$ and speed $v$ measured in a single inertial frame. The speed of light in vacuum is denoted by $c$. Define the dimensionless speed ratio and Lorentz factor by
\begin{equation}
  \beta=\frac{v}{c},
  \qquad
  \gamma=\frac{1}{\sqrt{1-\beta^2}}.
  \label{eq:velocity-lorentz-factor}
\end{equation}
For a massive particle, the physical range is $0\leq v<c$, or $0\leq\beta<1$.

The relativistic energy--momentum relation is
\begin{equation}
  E^2=p^2c^2+m^2c^4,
  \label{eq:energy-momentum-relation}
\end{equation}
where $E$ and momentum magnitude $p$ refer to the same inertial frame. Using $p=\gamma mv$ gives
\begin{align}
  E^2
  &=\gamma^2m^2v^2c^2+m^2c^4 \label{eq:energy-momentum-substitution} \\
  &=m^2c^4\left(\gamma^2\beta^2+1\right) \label{eq:energy-momentum-simplification}
   =\gamma^2m^2c^4,
\end{align}
because $\gamma^2\beta^2+1=\gamma^2$. Taking the positive root yields the exact total energy
\begin{equation}
  E=\gamma mc^2
   =\frac{mc^2}{\sqrt{1-v^2/c^2}}.
  \label{eq:exact-energy}
\end{equation}
Equation~\eqref{eq:exact-energy} is the starting point for the expansion.

At zero speed, the particle has rest energy
\begin{equation}
  E_0=mc^2.
  \label{eq:rest-energy}
\end{equation}
The kinetic energy in the chosen frame is defined as the excess of total energy over rest energy:
\begin{equation}
  K=E-E_0=(\gamma-1)mc^2.
  \label{eq:kinetic-definition}
\end{equation}
Thus $E$ is frame-dependent, while $E_0$ is set by the invariant mass. The subtraction in Eq.~\eqref{eq:kinetic-definition} isolates the energy associated with the particle's motion relative to the selected frame.

\section{Choosing a dimensionless expansion variable}

The speed occurs in Eq.~\eqref{eq:exact-energy} as the ratio $v^2/c^2$. Set
\begin{equation}
  x=\beta^2=\frac{v^2}{c^2},
  \qquad
  f(x)=(1-x)^{-1/2}.
  \label{eq:dimensionless-function}
\end{equation}
The total energy can then be written as $E=mc^2f(x)$. This change of variable separates the dimensionless mathematical expansion from the dimensional factor $mc^2$.

For a function analytic around zero, its Taylor series is
\begin{equation}
  f(x)=\sum_{n=0}^{\infty}a_nx^n,
  \qquad
  a_n=\frac{f^{(n)}(0)}{n!}.
  \label{eq:taylor-definition}
\end{equation}
For the function in Eq.~\eqref{eq:dimensionless-function}, repeated differentiation gives
\begin{equation}
  f^{(n)}(x)=\left(\frac{1}{2}\right)_n(1-x)^{-n-1/2},
  \label{eq:binomial-derivative}
\end{equation}
where $(1/2)_n=(1/2)(3/2)(5/2)\cdots((2n-1)/2)$ is the rising factorial, with $(1/2)_0=1$. Evaluating at zero and dividing by $n!$ gives
\begin{equation}
  a_n=\frac{(1/2)_n}{n!}
     =\frac{(2n)!}{4^n(n!)^2}
     =\frac{\binom{2n}{n}}{4^n}.
  \label{eq:general-coefficient}
\end{equation}
The factorial in the Taylor denominator matters: derivatives at zero are not themselves the coefficients.

The first derivatives and Taylor coefficients illustrate the distinction:
\begin{align*}
  f(0)&=1, & a_0&=1, \\
  f'(0)&=\frac{1}{2}, & a_1&=\frac{1}{2}, \\
  f''(0)&=\frac{3}{4}, & a_2&=\frac{3}{8}, \\
  f'''(0)&=\frac{15}{8}, & a_3&=\frac{5}{16}.
\end{align*}
The Taylor series is therefore
\begin{equation}
  (1-x)^{-1/2}
  =1+\frac{1}{2}x+\frac{3}{8}x^2+\frac{5}{16}x^3
   +\frac{35}{128}x^4+\cdots.
  \label{eq:dimensionless-series}
\end{equation}
This is the binomial expansion with exponent $-1/2$ and argument $-x$ \cite{DLMF}. Its radius of convergence is one, so Eq.~\eqref{eq:dimensionless-series} represents the function for $|x|<1$.

\section{Restoring units and separating the energy terms}

Substitute $x=v^2/c^2$ into Eq.~\eqref{eq:dimensionless-series} and multiply by $mc^2$. The result is
\begin{align}
  E
  &=mc^2\left[
    1+\frac{1}{2}\frac{v^2}{c^2}
      +\frac{3}{8}\frac{v^4}{c^4}
      +\frac{5}{16}\frac{v^6}{c^6}
      +\frac{35}{128}\frac{v^8}{c^8}+\cdots
    \right] \\
  &=mc^2+\frac{1}{2}mv^2
      +\frac{3}{8}\frac{mv^4}{c^2}
      +\frac{5}{16}\frac{mv^6}{c^4}
      +\frac{35}{128}\frac{mv^8}{c^6}+\cdots.
  \label{eq:energy-series}
\end{align}
Every term has units of energy. For example, $mv^4/c^2$ has units of mass times speed squared, as required. Only even powers of speed appear because the exact energy depends on $v^2$.

The $n=0$ term is the rest energy $mc^2$. Removing it according to Eq.~\eqref{eq:kinetic-definition} gives the kinetic-energy series:
\begin{equation}
  K=\frac{1}{2}mv^2
      +\frac{3}{8}\frac{mv^4}{c^2}
      +\frac{5}{16}\frac{mv^6}{c^4}
      +\frac{35}{128}\frac{mv^8}{c^6}+\cdots.
  \label{eq:kinetic-series}
\end{equation}
Equivalently, the exact infinite-series forms are
\begin{align}
  E&=mc^2\sum_{n=0}^{\infty}
       \frac{\binom{2n}{n}}{4^n}\left(\frac{v}{c}\right)^{2n},
       \label{eq:energy-sum}\\
  K&=mc^2\sum_{n=1}^{\infty}
       \frac{\binom{2n}{n}}{4^n}\left(\frac{v}{c}\right)^{2n}.
       \label{eq:kinetic-sum}
\end{align}
Both sums converge for $v<c$.

Equation~\eqref{eq:kinetic-series} makes the Newtonian limit explicit. When $v/c$ is small, the first term dominates:
\begin{equation}
  K\simeq\frac{1}{2}mv^2.
  \label{eq:low-speed-kinetic-energy}
\end{equation}
The first relativistic correction is $3mv^4/(8c^2)$. Relative to the Newtonian term, it is $3(v/c)^2/4$. Thus the leading fractional correction to the Newtonian kinetic energy is small when $v\ll c$ and grows with the square of the speed ratio.

\section{Convergence and finite truncations}

The Taylor series has radius of convergence one in $x$. Since $x=(v/c)^2$, every physical speed $0\leq v<c$ lies inside that domain. At $v=c$, the exact Lorentz factor diverges and the series does not converge; the boundary is not a valid point at which to use the expansion. This domain statement concerns the infinite series. It does not by itself guarantee that a short partial sum is accurate at every speed within the domain.

To make that distinction precise, define the approximation through order $v^{2N}$ by
\begin{equation}
  E_N=mc^2\sum_{n=0}^{N}
       \frac{\binom{2n}{n}}{4^n}\left(\frac{v}{c}\right)^{2n}.
  \label{eq:truncation}
\end{equation}
For example, the approximation through $v^4$ retains terms through $n=2$:
\begin{equation}
  \frac{E_2}{E_0}=1+\frac{1}{2}\beta^2+\frac{3}{8}\beta^4.
  \label{eq:through-v4}
\end{equation}
All coefficients in Eq.~\eqref{eq:energy-sum} are positive. Consequently, for $0<v<c$, any finite partial sum is below the exact total energy, and omitted terms make a positive contribution. As $v$ approaches $c$, more terms are needed to approximate the rapidly increasing exact energy.

Table~\ref{tab:comparison} compares the exact energy with Eq.~\eqref{eq:through-v4}. The relative shortfall is defined with respect to the exact total energy:
\begin{equation}
  \delta_2=100\,\frac{E-E_2}{E}\ \text{percent}.
  \label{eq:relative-shortfall}
\end{equation}
Values are rounded. The comparison is for total energy, including rest energy; a relative error computed for kinetic energy alone would use a different denominator.

\begin{table}[ht]
  \centering
  \caption{Exact total energy and the approximation through $v^4$. Energies are normalized by $E_0=mc^2$.}
  \label{tab:comparison}
  \begin{tabular}{cccc}
    \toprule
    $\beta=v/c$ & Exact $E/E_0$ & Through $v^4$, $E_2/E_0$ & Relative shortfall \\
    \midrule
    0.1 & 1.0050378 & 1.0050375 & 0.00003\% \\
    0.5 & 1.1547005 & 1.1484375 & 0.54\% \\
    0.9 & 2.2941573 & 1.6510375 & 28.03\% \\
    \bottomrule
  \end{tabular}
\end{table}

At $v=0.1c$, the short sum is an excellent approximation to the total energy. At $v=0.5c$, the error is still below one percent for this truncation. At $v=0.9c$, the same two-term correction beyond rest energy misses a substantial part of the exact total energy. The exact infinite series remains valid at all three speeds, but the usefulness of a finite truncation depends on both the speed and the number of terms retained.

\section{Physical interpretation and limits}

Several distinctions help prevent overinterpreting the series.

\paragraph{Total energy and kinetic energy.}
The constant term in Eq.~\eqref{eq:energy-series} is not kinetic energy. It is rest energy, present even at $v=0$. Kinetic energy is obtained only after subtracting $mc^2$. In particular, the low-speed approximation applies to $K$, whose leading term is $mv^2/2$; it does not mean that the total energy is approximately just $mv^2/2$.

\paragraph{Exact series and approximation.}
Within its convergence domain, the full infinite sum is another representation of the exact relativistic expression. A finite sum is an approximation. The first statement is about convergence; the second is about numerical accuracy. Near the speed of light, one can remain inside the convergence domain while a short truncation is still poor.

\paragraph{Continuous corrections, not energy levels.}
The powers of $v/c$ label terms in a mathematical expansion. They do not imply that a particle occupies discrete energy levels. The series describes how a continuous classical energy function varies with speed.

\paragraph{Scope of the model.}
This calculation concerns a free massive particle and its translational energy in one inertial frame. It does not include potential energy, interactions, radiation, or gravitational effects. The invariant mass $m$ is held fixed. The expansion also assumes the standard special-relativistic expression in Eq.~\eqref{eq:exact-energy}; it does not derive that expression from more basic postulates. Einstein's 1905 paper, Section 10, discusses the kinetic energy of an electron in its own historical derivation \cite{Einstein1905}, which is related context but has a different starting point and purpose.

\section{Conclusion}

Expanding the exact relativistic total energy around zero speed gives a transparent hierarchy of terms. The constant term is the rest energy $mc^2$. After it is subtracted, the leading kinetic term is the Newtonian expression $mv^2/2$. Higher even powers of speed provide relativistic corrections, with the first correction proportional to $mv^4/c^2$.

The dimensionless variable $x=(v/c)^2$ exposes the binomial structure and shows that the infinite series converges for $|v|<c$. For physical speeds this is the full subluminal range, but convergence does not make every short truncation accurate near $c$. The central lesson is therefore twofold: Newtonian kinetic energy emerges as the low-speed limit, and the exact relativistic expression should be retained when a finite low-order approximation is insufficient.

\appendix
\section{A recurrence for the coefficients}

The general coefficient in Eq.~\eqref{eq:general-coefficient} is $a_n=\binom{2n}{n}/4^n$. Its adjacent-term ratio and recurrence are:
\begin{equation}
  \frac{a_{n+1}}{a_n}=\frac{2n+1}{2n+2},
  \qquad
  a_{n+1}=\frac{2n+1}{2n+2}a_n,
  \qquad a_0=1.
  \label{eq:coefficient-recurrence}
\end{equation}
For example, $a_1=1/2$, $a_2=(3/4)(1/2)=3/8$, and $a_3=(5/6)(3/8)=5/16$. This recurrence is convenient for computing further terms without repeatedly evaluating factorials.

\begin{thebibliography}{9}
\small
\setstretch{1.0}
\raggedright
\urlstyle{same}

\bibitem{carousel}
Light-Meter Physics,
\emph{Relativistic Energy / Taylor Series}, instructional carousel, seven slides,
undated source PDF supplied with this manuscript.

\bibitem{DLMF}
F. W. J. Olver, A. B. Olde Daalhuis, D. W. Lozier, B. I. Schneider,
R. F. Boisvert, C. W. Clark, B. R. Miller, B. V. Saunders, H. S. Cohl,
and M. A. McClain, eds.,
\emph{NIST Digital Library of Mathematical Functions}, Section 4.6(ii),
Eq.~4.6.7, National Institute of Standards and Technology,
\url{https://dlmf.nist.gov/4.6#E7}.

\bibitem{Einstein1905}
A. Einstein, ``On the Electrodynamics of Moving Bodies,''
\emph{Annalen der Physik} \textbf{17}, 891--921 (1905), Section 10.
English translation in H. A. Lorentz et al.,
\emph{The Principle of Relativity} (Methuen, London, 1923).

\bibitem{EnergyMomentum}
Light-Meter Physics,
\emph{Relativistic Energy as a Function of Momentum: Low- and High-Momentum Expansions},
companion manuscript in the Papers series, draft.

\bibitem{ProperTime}
Light-Meter Physics,
\emph{Proper Time and Kinetic Energy in Special Relativity: A Rate Expansion and an Exact Identity},
companion manuscript in the Papers series, draft.

\end{thebibliography}

\end{document}